\documentclass[9pt,a4paper]{article}
\usepackage[utf8]{inputenc}
\usepackage[T1]{fontenc}
\usepackage{lmodern}
\usepackage{amsmath,amssymb,bm}
\usepackage[margin=1.2cm]{geometry}
\usepackage{xcolor}
\usepackage{titlesec}
\usepackage{enumitem}
\usepackage{booktabs}
\usepackage{microtype}
\usepackage{hyperref}
\usepackage{tcolorbox}

\definecolor{cnamred}{RGB}{196, 18, 45}
\definecolor{cnamdark}{RGB}{120, 10, 28}
\definecolor{boxbg}{RGB}{253, 248, 249}
\definecolor{graytext}{RGB}{80, 80, 80}

\titleformat{\section}{\normalsize\bfseries\color{cnamred}}{\thesection}{1em}{}[\titlerule]
\newtcolorbox{slidebox}[2][]{colback=boxbg,colframe=cnamred!75,fonttitle=\bfseries\footnotesize,title={#2},top=2pt,bottom=2pt,left=5pt,right=5pt,#1}

\setlist[itemize]{leftmargin=1.2em, itemsep=1pt, topsep=1pt}

\title{\textbf{\color{cnamred} Master 2 Oral Defense: Slide-by-Slide Spoken Script}\\[0.1cm]\small\color{graytext} Real-Time Physics on the Web: IGA \& ROM for Parameterized Structural Dynamics}
\author{\textbf{Mumen Wehbe} \quad$\bullet$\quad Supervisor: \textbf{Dr. Christophe Hoareau} \quad$\bullet$\quad LMSSC / Cnam Paris}
\date{Target Time: \textbf{28:00} / 30:00 Max \quad$\bullet$\quad Short, Punchy Sentences for Every Single Slide}

\begin{document}
\maketitle

\vspace{-0.5cm}
\begin{center}
\scriptsize
\begin{tabular}{lcccc}
\toprule
\textbf{Chapter} & \textbf{Pages} & \textbf{Clock} & \textbf{Duration} & \textbf{Key Action / Demonstrations} \\
\midrule
Preamble \& Overview & 1 -- 4 & 00:00 -- 02:00 & 2m 00s & Title, LMSSC team, agenda overview \\
Ch 1: Intro \& Motivation & 5 -- 7 & 02:00 -- 04:15 & 2m 15s & Dual scope, vector/audio compression parallels \\
Ch 2: Theory \& Pullback & 8 -- 13 & 04:15 -- 08:30 & 4m 15s & Parameters $(r, x_c)$, moving domain, $\hat{\Omega}$ pullback \\
Ch 3: Isogeometric Analysis & 14 -- 25 & 08:30 -- 15:30 & 7m 00s & \textbf{Live Demo 1} (Page 19: B-Spline curve), $k$-refinement \\
Ch 4: Full Order Model (FOM) & 26 -- 32 & 15:30 -- 19:45 & 4m 15s & Newmark-$\beta$, Von Mises stress, $<10^{-12}$ validation \\
Ch 5: ROM \& Hyper-Reduction & 33 -- 41 & 19:45 -- 26:15 & 6m 30s & POD, ECSW (28 elements), \textbf{Grand Live Demo 2} (Page 40) \\
Ch 6: Conclusion \& Perspectives & 42 -- 45 & 26:15 -- 28:00 & 1m 45s & Summary of 4 pillars, future axes, open discussion \\
\bottomrule
\end{tabular}
\end{center}

\vspace{-0.1cm}

% ── PREAMBLE ──────────────────────────────────────────────────────────────────
\section*{Preamble (Slides 1 -- 4)}

\begin{slidebox}{Page 1 | Title Slide \hfill [00:00 -- 00:35 | 35s]}
\textit{[Stand center-stage, make eye contact with jury.]}
\begin{itemize}
  \item ``Good morning, Mr. President, esteemed members of the jury.''
  \item ``Welcome to my Master 2 defense in Computational Mechanics.''
  \item ``My topic is: \textit{Real-Time Physics on the Web: IGA and Reduced Order Modeling for Structural Dynamics}.''
  \item ``This research was supervised by Dr. Christophe Hoareau here at Cnam Paris.''
\end{itemize}
\textbf{\color{cnamred}[CLICK NEXT]}
\end{slidebox}

\begin{slidebox}{Page 2 | Host Institution \& Research Laboratory \hfill [00:35 -- 01:05 | 30s]}
\textit{[Point to the Saint-Martin portal photo.]}
\begin{itemize}
  \item ``This research was hosted at LMSSC, the Structural Mechanics and Coupled Systems Laboratory.''
  \item ``The laboratory specializes in structural dynamics, vibration mitigation, and fluid-structure interaction.''
  \item ``This work is part of the DYSCOM research group, focusing on interactive numerical methods.''
\end{itemize}
\textbf{\color{cnamred}[CLICK NEXT]}
\end{slidebox}

\begin{slidebox}{Page 3 | Supervision \& Research Team \hfill [01:05 -- 01:30 | 25s]}
\textit{[Acknowledge advisor with open hand.]}
\begin{itemize}
  \item ``I want to warmly thank my supervisor, Dr. Christophe Hoareau.''
  \item ``His scientific guidance and continuous support were invaluable throughout this project.''
  \item ``I also extend my sincere thanks to the professors and research staff at LMSSC.''
\end{itemize}
\textbf{\color{cnamred}[CLICK NEXT]}
\end{slidebox}

\begin{slidebox}{Page 4 | Lecture Overview (Table of Contents) \hfill [01:30 -- 02:00 | 30s]}
\textit{[Sweep through the 6 chapters on screen.]}
\begin{itemize}
  \item ``Today's talk is organized into six parts.''
  \item ``We start with the motivation and continuum mechanics theory.''
  \item ``Then we detail the IGA discretization and our Full Order Model solver.''
  \item ``Next, we present Model Order Reduction with POD and ECSW hyper-reduction.''
  \item ``Finally, I will present a live web demonstration and conclude.''
\end{itemize}
\textbf{\color{cnamred}[CLICK NEXT]}
\end{slidebox}

% ── CHAPTER 1 ─────────────────────────────────────────────────────────────────
\section*{Chapter 1: Introduction \& State of the Art (Slides 5 -- 7)}

\begin{slidebox}{Page 5 | Chapter 1 Announcement \hfill [02:00 -- 02:10 | 10s]}
\begin{itemize}
  \item ``Let us begin with Chapter 1: Introduction and Motivation.''
\end{itemize}
\textbf{\color{cnamred}[CLICK NEXT]}
\end{slidebox}

\begin{slidebox}{Page 6 | Introduction \& Motivation \hfill [02:10 -- 03:15 | 1m 05s]}
\textit{[Point to left blocks; mention web screenshot on the right as the target goal.]}
\begin{itemize}
  \item ``Our project has two clear goals:''
  \item ``\textbf{Education}: an interactive physics visualizer directly in the browser with zero install.''
  \item ``Students can physically see wave propagation and vibrations live.''
  \item ``\textbf{Research}: solving the moving-geometry problem in real time.''
  \item ``Standard FEM is too slow and crashes on remeshing.''
  \item ``Our solution: exact CAD geometry with IGA, and extreme speed with ROM.''
\end{itemize}
\textbf{\color{cnamred}[CLICK NEXT]}
\end{slidebox}

\begin{slidebox}{Page 7 | Conceptual Parallels to Data Compression \hfill [03:15 -- 04:15 | 1m 00s]}
\textit{[Point to Hughes (left), Quarteroni (center), and Farhat (right).]}
\begin{itemize}
  \item ``We can understand our pipeline through direct parallels to data compression:''
  \item ``1. \textbf{IGA (Hughes, 2005)} is like vector graphics: CAD curves are exact, never pixelated into chords.''
  \item ``2. \textbf{ROM (Quarteroni, 2016)} is like MP3 audio compression: we retain only dominant dynamic modes.''
  \item ``3. \textbf{Hyper-Reduction (Farhat, 2014)} is like smart sensor placement: we compute on just a few elements.''
\end{itemize}
\textbf{\color{cnamred}[CLICK NEXT]}
\end{slidebox}

% ── CHAPTER 2 ─────────────────────────────────────────────────────────────────
\section*{Chapter 2: Theory: Problem Definition (Slides 8 -- 13)}

\begin{slidebox}{Page 8 | Chapter 2 Announcement \hfill [04:15 -- 04:25 | 10s]}
\begin{itemize}
  \item ``Now we move to Chapter 2: Theoretical Formulation and Continuum Mechanics.''
\end{itemize}
\textbf{\color{cnamred}[CLICK NEXT]}
\end{slidebox}

\begin{slidebox}{Page 9 | Geometry Configurations \hfill [04:25 -- 05:15 | 50s]}
\textit{[Point to cantilever beam schematic, indicating $r$ and $x_c$.]}
\begin{itemize}
  \item ``We study a 2D clamped cantilever beam with symmetric smooth notches.''
  \item ``The geometry depends on two parameters: notch depth $r$ and notch position $x_c$.''
  \item ``When these parameters change, the physical domain deforms.''
  \item ``In standard FEM, this requires remeshing for every single parameter change.''
\end{itemize}
\textbf{\color{cnamred}[CLICK NEXT]}
\end{slidebox}

\begin{slidebox}{Page 10 | Governing Equations \& Assumptions \hfill [05:15 -- 06:00 | 45s]}
\textit{[Point to linear momentum equation.]}
\begin{itemize}
  \item ``We assume linear elasticity, small deformations, and 2D plane stress.''
  \item ``The motion is governed by the elastodynamic momentum balance equation.''
  \item ``The left boundary is fully clamped, and a dynamic traction acts at the tip.''
\end{itemize}
\textbf{\color{cnamred}[CLICK NEXT]}
\end{slidebox}

\begin{slidebox}{Page 11 | Classical Weak Form on $\Omega(\boldsymbol{\mu})$ \hfill [06:00 -- 06:50 | 50s]}
\textit{[Highlight the moving domain subscript $\Omega(\boldsymbol{\mu})$.]}
\begin{itemize}
  \item ``Here is the classical weak form on the physical domain.''
  \item ``Notice the integration domain $\Omega(\boldsymbol{\mu})$ moves with the parameter.''
  \item ``This is the central bottleneck: snapshots belong to incompatible function spaces.''
  \item ``We cannot build a global reduced basis directly from these meshes.''
\end{itemize}
\textbf{\color{cnamred}[CLICK NEXT]}
\end{slidebox}

\begin{slidebox}{Page 12 | Pullback Mapping to Reference Configuration \hfill [06:50 -- 07:40 | 50s]}
\textit{[Trace mapping arrow from $\hat{\Omega}$ down to $\Omega(\boldsymbol{\mu})$.]}
\begin{itemize}
  \item ``To fix this, we map the moving domain to a fixed reference rectangle $\hat{\Omega}$.''
  \item ``This is the reference pullback mapping $\boldsymbol{\Psi}$.''
  \item ``All coordinate changes are captured by the mapping Jacobian matrix.''
  \item ``Now, all integrations take place on a static, parameter-independent space.''
\end{itemize}
\textbf{\color{cnamred}[CLICK NEXT]}
\end{slidebox}

\begin{slidebox}{Page 13 | Pulled-Back Weak Form on $\hat{\Omega}$ \hfill [07:40 -- 08:30 | 50s]}
\textit{[Point to Jacobian terms in stiffness and mass forms.]}
\begin{itemize}
  \item ``This gives our pulled-back weak form on $\hat{\Omega}$.''
  \item ``Notice that the integration bounds are completely fixed.''
  \item ``All parameter dependence is isolated inside the Jacobian mapping tensors.''
  \item ``Degrees of freedom now share a unified algebraic vector space, ideal for ROM.''
\end{itemize}
\textbf{\color{cnamred}[CLICK NEXT]}
\end{slidebox}

% ── CHAPTER 3 ─────────────────────────────────────────────────────────────────
\section*{Chapter 3: Discretization: Isogeometric Analysis (Slides 14 -- 25)}

\begin{slidebox}{Page 14 | Chapter 3 Announcement \hfill [08:30 -- 08:40 | 10s]}
\begin{itemize}
  \item ``Chapter 3 introduces the spatial discretization using Isogeometric Analysis.''
\end{itemize}
\textbf{\color{cnamred}[CLICK NEXT]}
\end{slidebox}

\begin{slidebox}{Page 15 | Curve Representations in Geometry \hfill [08:40 -- 09:20 | 40s]}
\textit{[Compare Explicit, Implicit, and Parametric plots.]}
\begin{itemize}
  \item ``Let us compare how curves are described in geometry.''
  \item ``Explicit curves cannot describe vertical lines or closed circles.''
  \item ``Implicit curves are computationally expensive to parameterize.''
  \item ``Parametric curves are the CAD standard: flexible, robust, and exact.''
\end{itemize}
\textbf{\color{cnamred}[CLICK NEXT]}
\end{slidebox}

\begin{slidebox}{Page 16 | Why IGA? \hfill [09:20 -- 10:05 | 45s]}
\textit{[Point to the three vertical blocks.]}
\begin{itemize}
  \item ``IGA unifies CAD geometry and numerical simulation:''
  \item ``\textbf{CAD Splines}: uses native spline functions, bypassing separate meshing steps.''
  \item ``\textbf{Exact Geometry}: preserves curved boundaries exactly at all scales without error.''
  \item ``\textbf{High Regularity}: $C^{p-1}$ continuity provides smooth stresses and higher accuracy per DOF.''
\end{itemize}
\textbf{\color{cnamred}[CLICK NEXT]}
\end{slidebox}

\begin{slidebox}{Page 17 | The Core Ingredients of IGA \hfill [10:05 -- 10:50 | 45s]}
\textit{[Point to the 4 ingredient blocks.]}
\begin{itemize}
  \item ``The IGA framework relies on four ingredients:''
  \item ``1. A knot vector $\Xi$ partitioning the parametric domain into elements.''
  \item ``2. B-spline basis functions computed by the Cox-de Boor recurrence.''
  \item ``3. Weights that generalize B-splines to NURBS for exact conics.''
  \item ``4. Control points defining the physical shape in Cartesian space.''
\end{itemize}
\textbf{\color{cnamred}[CLICK NEXT]}
\end{slidebox}

\begin{slidebox}{Page 18 | B-Spline Basis Functions \hfill [10:50 -- 11:35 | 45s]}
\textit{[Show basis curves and Cox-de Boor formula.]}
\begin{itemize}
  \item ``B-spline basis functions have ideal mathematical properties:''
  \item ``They form a partition of unity and are strictly non-negative.''
  \item ``They have compact local support, ensuring sparse system matrices.''
  \item ``They provide $C^{p-1}$ inter-element continuity across internal knots.''
\end{itemize}
\textbf{\color{cnamred}[CLICK NEXT]}
\end{slidebox}

\begin{slidebox}{Page 19 | B-Spline Curve Concept \& LIVE DEMO 1 \hfill [11:35 -- 12:35 | 1m 00s]}
\textit{[Action: Click button, switch to browser for 20 seconds, drag control point live.]}
\begin{itemize}
  \item ``A B-spline curve is a linear sum of basis functions multiplied by control points.''
  \item \textit{[Live Interaction]}: ``Notice as I drag this control point: the curve deforms smoothly and locally.''
  \item ``The knot topology never breaks, and no remeshing is needed.''
  \item ``This makes splines exceptionally suited for geometric parameterization.''
\end{itemize}
\textbf{\color{cnamred}[CLICK NEXT]}
\end{slidebox}

\begin{slidebox}{Page 20 | 2D Bivariate Basis Functions \hfill [12:35 -- 13:15 | 40s]}
\textit{[Show 2D surface patch.]}
\begin{itemize}
  \item ``For 2D continua, we take tensor products of 1D B-splines.''
  \item ``Each basis function is a simple product: $R_{i,j}(\xi, \eta) = N_i(\xi) M_j(\eta)$.''
  \item ``Control points form a structured 2D bidirectional net.''
\end{itemize}
\textbf{\color{cnamred}[CLICK NEXT]}
\end{slidebox}

\begin{slidebox}{Page 21 | Parametric to Physical Domain Mapping \hfill [13:15 -- 13:55 | 40s]}
\textit{[Trace mapping from unit square to physical notched beam.]}
\begin{itemize}
  \item ``This geometric mapping transforms the unit square $[0,1]^2$ into our beam.''
  \item ``Parametric knot spans map directly into physical elements.''
  \item ``Crucially, this parent square $[0,1]^2$ is our invariant reference domain $\hat{\Omega}$.''
\end{itemize}
\textbf{\color{cnamred}[CLICK NEXT]}
\end{slidebox}

\begin{slidebox}{Page 22 | $h$-Refinement (Knot Insertion) \hfill [13:55 -- 14:30 | 35s]}
\textit{[Show knot insertion plot.]}
\begin{itemize}
  \item ``To refine the mesh, we insert new knot values.''
  \item ``This splits elements into smaller elements, just like $h$-refinement in FEM.''
  \item ``Most importantly: inserting knots does not change the physical geometry at all.''
\end{itemize}
\textbf{\color{cnamred}[CLICK NEXT]}
\end{slidebox}

\begin{slidebox}{Page 23 | $p$-Refinement (Degree Elevation) \hfill [14:30 -- 15:00 | 30s]}
\textit{[Show degree elevation plot.]}
\begin{itemize}
  \item ``We can also elevate the polynomial degree $p$.''
  \item ``This enriches the approximation space without adding new elements.''
  \item ``Once again, the exact CAD geometry is strictly preserved.''
\end{itemize}
\textbf{\color{cnamred}[CLICK NEXT]}
\end{slidebox}

\begin{slidebox}{Page 24 | $k$-Refinement \& Exact Geometry \hfill [15:00 -- 15:35 | 35s]}
\textit{[Point to comparison table: $h \to p$ vs. $k$-ref.]}
\begin{itemize}
  \item ``This brings us to $k$-refinement, unique to IGA.''
  \item ``If we insert knots first, then elevate degree ($h \to p$), continuity drops to $C^0$.''
  \item ``In $k$-refinement, we elevate degree first, then insert knots: keeping $C^{p-1}$ continuity.''
  \item ``This gives maximum accuracy per DOF and eliminates artificial stress jumps.''
\end{itemize}
\textbf{\color{cnamred}[CLICK NEXT]}
\end{slidebox}

\begin{slidebox}{Page 25 | Isoparametric Concept \& Field Approximation \hfill [15:35 -- 16:05 | 30s]}
\textit{[Point to equation of displacement field.]}
\begin{itemize}
  \item ``We follow the isoparametric concept.''
  \item ``The unknown displacement field $\mathbf{u}_h$ is interpolated using the exact same spline basis.''
  \item ``Geometry and physics share the identical mathematical representation.''
\end{itemize}
\textbf{\color{cnamred}[CLICK NEXT]}
\end{slidebox}

% ── CHAPTER 4 ─────────────────────────────────────────────────────────────────
\section*{Chapter 4: Simulation: Full Order Model (Slides 26 -- 32)}

\begin{slidebox}{Page 26 | Chapter 4 Announcement \hfill [16:05 -- 16:15 | 10s]}
\begin{itemize}
  \item ``Chapter 4 covers the Full Order Model dynamic simulation.''
\end{itemize}
\textbf{\color{cnamred}[CLICK NEXT]}
\end{slidebox}

\begin{slidebox}{Page 27 | From IGA to the Full Order Model (FOM) \hfill [16:15 -- 17:00 | 45s]}
\textit{[Point to semi-discrete matrix equation.]}
\begin{itemize}
  \item ``Applying Galerkin projection yields the semi-discrete equation of motion: $\mathbf{M}\ddot{\mathbf{u}} + \mathbf{K}\mathbf{u} = \mathbf{F}$.''
  \item ``System matrices are assembled using Gaussian quadrature over knot spans.''
  \item ``Our full-order model has $N = 6{,}564$ degrees of freedom.''
\end{itemize}
\textbf{\color{cnamred}[CLICK NEXT]}
\end{slidebox}

\begin{slidebox}{Page 28 | Simulation Path \& Solver Workflow \hfill [17:00 -- 17:45 | 45s]}
\textit{[Trace time integration flow.]}
\begin{itemize}
  \item ``We integrate in time using the implicit Newmark-$\beta$ method.''
  \item ``It is unconditionally stable with parameters $\beta = 1/4$ and $\gamma = 1/2$.''
  \item ``We apply a dynamic harmonic tip load and compute transient deflections.''
\end{itemize}
\textbf{\color{cnamred}[CLICK NEXT]}
\end{slidebox}

\begin{slidebox}{Page 29 | Model Displacement Field \hfill [17:45 -- 18:25 | 40s]}
\textit{[Show deformed beam mesh.]}
\begin{itemize}
  \item ``Here is the dynamic displacement response of the notched cantilever beam.''
  \item ``Bending waves propagate cleanly through the structure.''
  \item ``Clamped boundary conditions on the left wall are strictly enforced.''
\end{itemize}
\textbf{\color{cnamred}[CLICK NEXT]}
\end{slidebox}

\begin{slidebox}{Page 30 | Model Stress \& Von Mises Definition \hfill [18:25 -- 19:05 | 40s]}
\textit{[Point to stress concentration at the notch.]}
\begin{itemize}
  \item ``From displacement derivatives, we compute Cauchy stresses and Von Mises fields.''
  \item ``Stress smoothly concentrates around the notch roots.''
  \item ``Because IGA basis functions have $C^1$ continuity, stresses are smooth across elements.''
  \item ``No artificial post-processing smoothing is required.''
\end{itemize}
\textbf{\color{cnamred}[CLICK NEXT]}
\end{slidebox}

\begin{slidebox}{Page 31 | Physical vs. Reference Mapped Path: Validation \hfill [19:05 -- 19:45 | 40s]}
\textit{[Highlight $< 10^{-12}$ difference.]}
\begin{itemize}
  \item ``We validated our reference pullback solver against the physical solver.''
  \item ``The numerical difference between both paths is under $10^{-12}$, pure machine precision.''
  \item ``This proves the mathematical exactness of our pullback implementation.''
\end{itemize}
\textbf{\color{cnamred}[CLICK NEXT]}
\end{slidebox}

\begin{slidebox}{Page 32 | Mesh Convergence Study \hfill [19:45 -- 20:20 | 35s]}
\textit{[Show log-log convergence plot.]}
\begin{itemize}
  \item ``We conducted a rigorous mesh convergence study.''
  \item ``The asymptotic error matches the theoretical optimal rate $O(h^{p+1})$.''
  \item ``$k$-refined IGA converges significantly faster than classical quadratic FEM.''
\end{itemize}
\textbf{\color{cnamred}[CLICK NEXT]}
\end{slidebox}

% ── CHAPTER 5 ─────────────────────────────────────────────────────────────────
\section*{Chapter 5: Reduced Order Modeling \& Hyper-Reduction (Slides 33 -- 41)}

\begin{slidebox}{Page 33 | Chapter 5 Announcement \hfill [20:20 -- 20:30 | 10s]}
\begin{itemize}
  \item ``Now we enter Chapter 5: Model Order Reduction and Hyper-Reduction.''
\end{itemize}
\textbf{\color{cnamred}[CLICK NEXT]}
\end{slidebox}

\begin{slidebox}{Page 34 | Why ROM? (Motivation) \hfill [20:30 -- 21:15 | 45s]}
\textit{[Contrast 142 ms FOM vs. 16 ms budget.]}
\begin{itemize}
  \item ``Solving the Full Order Model takes 142 milliseconds per time step.''
  \item ``To run at 60 FPS in a web browser, we need each step under 16 milliseconds.''
  \item ``The FOM is nearly 10 times too slow for real-time web deployment.''
  \item ``This is why we must build a Reduced Order Model.''
\end{itemize}
\textbf{\color{cnamred}[CLICK NEXT]}
\end{slidebox}

\begin{slidebox}{Page 35 | What Do We Need to Achieve It? (The 5 Steps) \hfill [21:15 -- 21:55 | 40s]}
\textit{[Walk through the 5 steps.]}
\begin{itemize}
  \item ``We reach real-time performance through five structured steps:''
  \item ``1. Reference pullback to fix the domain.''
  \item ``2. Offline snapshot collection across parameters.''
  \item ``3. POD basis construction via SVD.''
  \item ``4. ECSW hyper-reduction for fast assembly.''
  \item ``5. Client-side real-time solver in WebGL.''
\end{itemize}
\textbf{\color{cnamred}[CLICK NEXT]}
\end{slidebox}

\begin{slidebox}{Page 36 | Step 1: Reference Pullback Mapping \hfill [21:55 -- 22:30 | 35s]}
\textit{[Reiterate the fixed reference domain role.]}
\begin{itemize}
  \item ``Step 1 is the foundation: all geometries map to the same reference domain $\hat{\Omega}$.''
  \item ``Every snapshot shares the exact same algebraic degree-of-freedom indices.''
  \item ``This allows direct snapshot aggregation across the entire parameter space.''
\end{itemize}
\textbf{\color{cnamred}[CLICK NEXT]}
\end{slidebox}

\begin{slidebox}{Page 37 | Step 2: Snapshot Collection (Data Generation) \hfill [22:30 -- 23:05 | 35s]}
\textit{[Point to snapshot matrix structure.]}
\begin{itemize}
  \item ``In Step 2, we simulate training cases across the parameter space $\mathcal{D}$.''
  \item ``Displacement fields at each time step are stored as columns of snapshot matrix $\mathbf{S}$.''
  \item ``This captures the dominant dynamic deformation modes.''
\end{itemize}
\textbf{\color{cnamred}[CLICK NEXT]}
\end{slidebox}

\begin{slidebox}{Page 38 | Step 3: SVD / POD Basis Construction \hfill [23:05 -- 23:45 | 40s]}
\textit{[Point to singular value decay plot.]}
\begin{itemize}
  \item ``In Step 3, we compute the Singular Value Decomposition: $\mathbf{S} = \mathbf{U} \boldsymbol{\Sigma} \mathbf{V}^T$.''
  \item ``Singular values decay exponentially.''
  \item ``Just $r = 15$ POD modes capture over 99.99\% of total kinetic and strain energy.''
  \item ``We truncate the basis from 6,564 DOFs down to just 15.''
\end{itemize}
\textbf{\color{cnamred}[CLICK NEXT]}
\end{slidebox}

\begin{slidebox}{Page 39 | Step 4: Galerkin Projection vs. ECSW Hyper-reduction \hfill [23:45 -- 24:35 | 50s]}
\textit{[Explain the non-affine bottleneck.]}
\begin{itemize}
  \item ``Standard Galerkin projection shrinks matrix size, but still integrates over all 3,200 elements.''
  \item ``Online assembly still takes 150 milliseconds: no real speedup!''
  \item ``To fix this, we apply ECSW hyper-reduction: selecting a sparse subset of elements.''
  \item ``Using non-negative least squares, ECSW reduces 3,200 elements to just \textbf{28 active elements}.''
  \item ``Positive weights guarantee numerical stability, symmetry, and energy conservation.''
\end{itemize}
\textbf{\color{cnamred}[CLICK NEXT]}
\end{slidebox}

\begin{slidebox}{Page 40 | Step 5: Real-Time Solver \& GRAND LIVE DEMO 2 \hfill [24:35 -- 25:45 | 1m 10s]}
\textit{[Point to benchmark table, then switch to browser for 35 seconds.]}
\begin{itemize}
  \item ``Here are the results: solve time drops from 142 milliseconds to \textbf{1.3 milliseconds}.''
  \item ``That is a \textbf{64.5 times speedup}, with relative error below 0.1\%!''
  \item \textit{[Switch to browser]}: ``Here is our live Digital Twin running in the browser.''
  \item ``When I move notch depth $r$, the geometry and mesh update instantly.''
  \item ``When I move notch position $x_c$, dynamic stresses concentrate live at 60 FPS.''
  \item ``The entire solve runs client-side in WebGL with zero cloud latency.''
  \item \textit{[Switch back to slides]}
\end{itemize}
\textbf{\color{cnamred}[CLICK NEXT]}
\end{slidebox}

\begin{slidebox}{Page 41 | Project Workflow: Completing the Puzzle \hfill [25:45 -- 26:15 | 30s]}
\textit{[Trace the 5 workflow boxes.]}
\begin{itemize}
  \item ``This slide reflects our engineering methodology from start to finish:''
  \item ``1. Defined the core problem and real-time targets; 2. Structured the milestones with a Gantt chart;''
  \item ``3. Gathered references and studied prior literature (Hughes, Quarteroni, Farhat);''
  \item ``4. Solved formulation hurdles with Dr.~Hoareau's mentorship and AI pairing; 5. Delivered the final 60 FPS web application.''
\end{itemize}
\textbf{\color{cnamred}[CLICK NEXT]}
\end{slidebox}

% ── CHAPTER 6 ─────────────────────────────────────────────────────────────────
\section*{Chapter 6: Conclusion \& Future Perspectives (Slides 42 -- 45)}

\begin{slidebox}{Page 42 | Chapter 6 Announcement \hfill [26:15 -- 26:25 | 10s]}
\begin{itemize}
  \item ``We now conclude with Chapter 6: Summary and Perspectives.''
\end{itemize}
\textbf{\color{cnamred}[CLICK NEXT]}
\end{slidebox}

\begin{slidebox}{Page 43 | Summary of Achievements \& Implemented Methods \hfill [26:25 -- 27:05 | 40s]}
\textit{[Recap the 4 achievements with pride.]}
\begin{itemize}
  \item ``In summary, this research achieved four major milestones:''
  \item ``1. Developed an exact-geometry IGA dynamic elastodynamics solver.''
  \item ``2. Resolved the moving-boundary issue using reference pullback mapping.''
  \item ``3. Implemented POD and ECSW hyper-reduction, achieving a 64.5 times speedup on 28 elements.''
  \item ``4. Built a real-time web Digital Twin running at 60 FPS on standard devices.''
\end{itemize}
\textbf{\color{cnamred}[CLICK NEXT]}
\end{slidebox}

\begin{slidebox}{Page 44 | Alternatives \& Future Perspectives \hfill [27:05 -- 27:45 | 40s]}
\textit{[Present future research vectors.]}
\begin{itemize}
  \item ``For future research, three direct paths open:''
  \item ``First, extending to large geometric deformations with Green-Lagrange strain.''
  \item ``Second, advancing from 2D plane stress to 3D solid structures using trivariate splines.''
  \item ``Third, exploring neural networks to predict reduced coordinates for rapid design sweeps.''
\end{itemize}
\textbf{\color{cnamred}[CLICK NEXT]}
\end{slidebox}

\begin{slidebox}{Page 45 | Thank You / Questions \& Discussion \hfill [27:45 -- 28:00 | 15s]}
\textit{[Stand upright, smile, make eye contact, bow slightly.]}
\begin{itemize}
  \item ``Thank you very much for your kind attention.''
  \item ``I am now pleased to answer any questions from the jury.''
\end{itemize}
\textbf{\color{cnamred}[FINISHED AT 28:00 -- 2:00 SAFETY BUFFER BEFORE 30:00 LIMIT]}
\end{slidebox}

\end{document}
